%!TEX root = ./../../springer-mono-format/book.tex
%!TEX root = ./../../thesis-trier.tex

Here, we show the following:

\begin{enumerate}
	\item \ST, \DST, and \DFAINT are \NP-complete for solvable input automata
	(\cref{thm:solvable}).
	
	\item \SFS is \NP-complete for solvable automata (\cref{thm:SF-solvable}).
	
	\item \DFAST, \DFASS, \DFAS, and \SIS are \LOGSPACE-complete
	for solvable automata even for a unary alphabet (\cref{thm:SYNC-solv}).
	
	\item \REACHF and \REACH are \NP-complete for solvable automata
	in (\cref{prelim:compl:thm:reach-NP,prelim:compl:thm:reachf-NP}).
\end{enumerate}
For self-containment, we give a proof of \NP-hardness of \DFAINT 
for unary permutation automata. The proof is similar to the one given in \cite[Lem.~27]{MartensN07},
but uses automata instead of regular expressions.

\begin{proposition}
	\label{compl:prop:dfaint_unary}
	\DFAINT is \NP-hard for unary permutation automata.
\end{proposition}
\begin{proof}	
	We give a reduction from {\sc 3SAT}, which 
	is \NP-complete by \cref{prelim:thm:THREESAT}.
	Let $C_1, \ldots, C_m \subseteq \{x_1, \overline x_1, \ldots, x_n, \overline x_n\}$ 
	be the clauses of an {\sc 3SAT}-instance over the variables $x_1, \ldots, x_n$.
	We can assume every variable appears, as a positive or negative literal, in some clause.
	
	\medskip 
	\noindent\underline{Construction:} 
	Let $p_1 < p_2 < \cdots < p_n$
	be the first $n$ prime numbers.
	For each clause $C_j$, $j \in \{1,\ldots,m\}$,
	we construct 
	$\automat{A_j}{Q_j}{\{a\}}{\delta_j}{q_j}{F_j}$
	with
	\begin{align*}
		Q_j & = \{0,1,\ldots, N_j - 1\}, \\
		q_j & = 0, \\
		F_j & = \{ q \in Q_j \mid \exists i \in \{1,\ldots,n\} : 
		( x_i \in C_j \land q \equiv_{p_i} 0 ) \lor ( \overline{x}_i \in C_j \land q \not\equiv_{p_i} 0 ) \}, \\
		%F_j & = \{ q \in Q_j \mid \exists i \in \{1,\ldots,n\} : 
		%( x_i \in C_j \land q \equiv 0 \pmod{p_i} ) \\ 
		%& \qquad\qquad\qquad\qquad\qquad\qquad\qquad \lor ( \overline{x}_i \in C_j \land q \not\equiv 0 \pmod{p_i}) \}, \\
		\applytrans{\delta_j}{q}{a} & = (q + 1) \bmod{ N_j }
	\end{align*}
	where $N_j = \left(\prod_{\substack{ i \in \{1,\ldots,n\} \\ x_i \in C_j \lor \overline{x}_i \in C_j}} p_i \right)$,
	i.e., we multiply those primes that correspond
	to the variables in~$C_j$.
	
	By \cref{prelim:cor:PNT} 
	there exists $N \ge 1$ such that $p_n < n^2$
	for all $n \ge N$.
	We can assume $p_n < n^2$ in the following.
	Since  $|C_j| \le 3$
	we get  $|Q_j| \le n^6$.
	Then, as detecting primality for unary numbers
	is in~\LOGSPACE~\cite[Thm.~7]{DBLP:journals/tcs/HartmanisB76},
	the automata $\mathcal A_j$ can be constructed 
	in deterministic log-space.	
	
	\Cref{prelim:compl:fig:intersection-unary-construction} 
	shows the construction for the clauses $C_1 = \{ x_1, x_2\}$, $C_2 = \{ \overline{x_1}, x_3\}$,
	and $C_3 = \{ x_2, \overline{x_3}\}$ (shown in their logical interpretation
	in the figure).
	If $x_1 = 1$, then $x_3 = 1$ to fulfill $C_2$
	and thus $x_2 = 1$ to fulfill $C_3$. For $x_1 = 0$,
	we must have $x_2 = 1$ to fulfill $C_1$
	and $x_3$ can be arbitrary. Thus we find 
	three satisfying assignments that correspond to 
	the words $a^{30}$, $a^3$, and $a^{15}$.	
	
	\begin{figure}[ht]
		\centering	
		
		\begin{tikzpicture}[>=Stealth] %[>=Stealth[length=3pt]] %
			
			% ---------- Stil & Radius ----------
			\def\radius{11em}
			\def\radiusA{0.3*\radius}
			\def\radiusB{0.44*\radius}
			\def\radiusC{0.67*\radius}
			\def\ylabel{\radiusC-2em}
			
			\tikzset{
				smallstate/.style={
					draw,
					circle,
					minimum size=15pt,
					inner sep=0pt,
					font=\small
				}
			}
			
			% ---------- Automat mit 6 Zuständen ----------
			\begin{scope}
				\foreach \i in {0,...,5} {
					% teilbar durch 2 ODER 3
					\pgfmathtruncatemacro{\isfinal}{(mod(\i,2)==0) || (mod(\i,3)==0)}
					\ifnum\isfinal=1
					\node[smallstate,accepting] (q\i) at ({360/6*\i}:\radiusA) {$\i$};
					\else
					\node[smallstate] (q\i) at ({360/6*\i}:\radiusA) {$\i$};
					\fi
				}
				\foreach \i in {0,...,5} {
					\pgfmathtruncatemacro{\j}{mod(\i+1,6)}
					\draw[->] (q\i) edge (q\j);
				}
				\node at (0,-\ylabel) {$\mathcal A_1$ for $C_1 = x_1 \lor x_2$};
			\end{scope}
			
			% ---------- Automat mit 10 Zuständen ----------
			\begin{scope}[xshift=12em]
				\foreach \i in {0,...,9} {
					% ungerade UND durch 3 teilbar
					\pgfmathtruncatemacro{\isfinal}{(mod(\i,2)==1) || (mod(\i,5)==0)}
					\ifnum\isfinal=1
					\node[smallstate,accepting] (p\i) at ({360/10*\i}:\radiusB) {$\i$};
					\else
					\node[smallstate] (p\i) at ({360/10*\i}:\radiusB) {$\i$};
					\fi
				}
				\foreach \i in {0,...,9} {
					\pgfmathtruncatemacro{\j}{mod(\i+1,10)}
					\draw[->] (p\i) edge (p\j);
				}
				\node at (0,-\ylabel) {$\mathcal A_2$ for $C_2 = \overline{x_1} \lor x_3$};
			\end{scope}
			
			% ---------- Automat mit 15 Zuständen ----------
			\begin{scope}[xshift=28em]
				\foreach \i in {0,...,14} {
					% gerade UND NICHT durch 5 teilbar
					\pgfmathtruncatemacro{\isfinal}{(mod(\i,2)==0) || (mod(\i,5)!=0)}
					\ifnum\isfinal=1
					\node[smallstate,accepting] (r\i) at ({360/15*\i}:\radiusC) {${\i}$};
					\else
					\node[smallstate] (r\i) at ({360/15*\i}:\radiusC) {${\i}$};
					\fi
				}
				\foreach \i in {0,...,14} {
					\pgfmathtruncatemacro{\j}{mod(\i+1,15)}
					\draw[->] (r\i) edge (r\j);
				}
				\node at (0,-\ylabel) {$\mathcal A_3$ for $C_3 = x_2 \lor \overline{x_3}$};
			\end{scope}	
		\end{tikzpicture}
		\caption[Example for reduction of \THREESAT to \DFAINT\ on unary automata]{Example for the proof of \cref{compl:prop:dfaint_unary}. States with label~$0$ are start states. Final states are indicated by double circles. The label~$a$ at the transitions is omitted.}
		\label{prelim:compl:fig:intersection-unary-construction}
	\end{figure} 
		
	\medskip
	\noindent\underline{Properties:} We have the following two properties:
	\begin{enumerate}
		\item \label{prelim:compl:INT-unary:1} $\delta_j(q_j, a^r) \equiv_{N_j} r$ for each $r \ge 0$.
		
		\item \label{prelim:compl:INT-unary:2} $(b_1, \ldots, b_n)$ is a satisfying assignment
		if and only if $a^{\prod_{i=1}^n p_i^{b_i}} \in \bigcap_{j=1}^m L(\mathcal A_j)$.
	\end{enumerate}	
	Property~\ref{prelim:compl:INT-unary:1} follows 
	from the construction of $\mathcal A_j$ whose automaton graph forms a single cycle.	
	Property~\ref{prelim:compl:INT-unary:2} follows since for $(b_1, \ldots, b_n) \in \{0,1\}^n$,
	$l = \prod_{i=1}^n p_i^{b_i}$,  $j \in \{1,\ldots,m\}$,
	and $s_j = \delta_j(q_j, a^l)$ we have $s_j \equiv l \pmod{N_j}$ and then
	\[
	\begin{array}{rlr}
		&\exists i \in I \colon ( ( x_i \in C_j \land b_i = 1 ) \lor ( \overline{x}_i \in C_j \land b_i = 0 ) ) \\
		& \Leftrightarrow \exists i \in I \colon ( ( x_i \in C_j \land l \equiv_{p_i} 0 ) \lor ( \overline{x}_i \in C_j \land l \not\equiv_{p_i} 0  ) )  & \text{(Def. $l$)} \\
		& \Leftrightarrow  \exists i \in I \colon( ( x_i \in C_j \land s_j \equiv_{p_i} 0 ) \lor ( \overline{x}_i \in C_j \land s_j \not\equiv_{p_i} 0 ) ) & \text{(Eqn.~\eqref{prelim:eqn:ab_divisor}, Property~(1))} \\
		& \Leftrightarrow  s_j \in F_j & \text{(Def. $F_j$)}
	\end{array}
	\]
	where $I = \{1,\ldots,n\}$.
	
	\medskip 
	\noindent\underline{Verification:} By Property~\ref{prelim:compl:INT-unary:2},
	if there exists a satisfying assignment, then we have $\bigcap_{j=1}^m L(\mathcal A_j) \ne \emptyset$.
	Conversely, suppose $a^l \in \bigcap_{j=1}^m L(\mathcal A_j)$
	for some $l \ge 0$. Define an assignment $(b_1, \ldots, b_n)$
	to the variables by setting, for $i \in \{1,\ldots,n\}$,
	\[
	b_i = \begin{cases}
		1 & \text{if } p_i \text{ divides } l, \\ 
		0 & \text{if } p_i \text{ does not divide } l. 
	\end{cases}	
	\] 	
	Let $l' = \prod_{i=1}^n p_i^{b_i}$.
	Then $l' \equiv_{p_i} 0$ if and only if $l \equiv_{p_i} 0$ for $i \in \{1,\ldots,n\}$.
	So $a^{l'} \in L(\mathcal A_j)$ if and only if $a^l \in L(\mathcal A_j)$ for $j \in \{1,\ldots,m\}$.
	Hence, by Property~\ref{prelim:compl:INT-unary:2}, we obtain that $(b_1, \ldots, b_n)$
	is a satisfying assignment.
\end{proof}

\begin{theorem}\label{thm:solvable} \label{compl:thm:comm_ST}
	\ST, \DST, and \DFAINT are \NP-complete for solvable input automata,
	where \NP-hardness holds
	even for the following classes of input automata:
	\begin{enumerate}
		\item binary partially ordered input automata,
		
		\item ternary partially ordered monotonic automata,
		
		\item commutative partially ordered  automata over an unbounded alphabet,
		
		\item unary permutation automata.
	\end{enumerate}
\end{theorem}
\begin{proof}
	That \DFAINT is \NP-complete for the mentioned
    classes follows from results
    in the literature mentioned in \cref{compl:sec:subclasses:known-results}, page~\pageref{prelim:classes:intro:INT}.
    By \cref{compl:prop:reductions:item:DFAINT_DST,compl:prop:reductions:item:DST_ST}
    from \cref{compl:prop:reductions} 
    we obtain \NP-hardness for all three problems,
    and by these statements it is sufficient to establish \NP-membership
    for \ST, which we do in the following.
    
	Let $\Sigma = \{a_1, \ldots, a_k\}$ and 
	$\semiautomat{A}{Q}{\Sigma}{\delta}$ be a solvable
	input automaton to \ST with $n$ states and input sets $S, T \subseteq Q$.
	If there exists a word mapping~$S$ into~$T$,
	by Proposition~\ref{prop:S-T-bound} there exists a word
	of the form
	\[ 
	v = v_0 x_1 \cdots v_{r-1} x_r v_r
	\]
	for some $r \le n^2$
	where $x_1, \ldots, x_r \in \Sigma$ and $v_i = a_1^{r_{i,1}} \cdots a_k^{r_{i,k}}$
	for numbers $0 \le r_{i,j} < 2^n$
	with $i \in \{0,1,\ldots,r\}$ and $j \in \{1,\ldots,k\}$
	mapping~$S$ into~$T$.
	By using a binary encoding for the exponents $0 \le r_{i,j} < 2^n$,
	a word of the above form can be stored using 
	$O( |\Sigma|n \cdot (n^2 + 1) + \log(|\Sigma|)n^2)$
	%$(|\Sigma|+1) n^3 + |\Sigma|n$ 
	bits where each of the $n^2+1$ words $v_i$, $i \in \{0,1,\ldots,r\}$,
	takes up $|\Sigma|n$ bits and each letter $x_i$, $i \in \{1,\ldots,r\}$,
	takes $\log( |\Sigma| )$ bits.
	So, using this encoding, it can be guessed in polynomial time.
	Next, the transition matrix
	of $a_j^{r_{i,j}}$ for $0 \le r_{i,j} < 2^n$
	can be computed in polynomial time by a square-and-multiply algorithm (also known as 
	exponentiation by squaring) 
	using the transition matrix of~$a_j$.
	Hence, the transition matrix of~$v$
	can be computed in polynomial time by multiplying
	the transition matrices of $v_0, x_1, \ldots, v_{r-1}, x_r$, and $v_r$ in this order.
	Then~$v$ maps~$S$ into~$T$
	if and only if each row in the transition matrix of~$v$ corresponding to a state in~$S$
	has precisely one~$1$ at a column corresponding to a state in~$T$.
	So, the above algorithm shows containment in \NP. 
\end{proof}

\begin{theorem}\label{thm:SF-solvable}
	\SFS is \NP-complete for solvable input automata even for binary partially ordered automata.
\end{theorem} 
\begin{proof}
	Ryzhikov~\cite[Thm.~4]{Ryzhikov19a} proved the problem to be \NP-complete
	for partially ordered binary automata, see also page~\pageref{prelim:classes:intro:SFS} 
	in \cref{compl:sec:subclasses:known-results}.
	
	For \NP-membership, we can guess a state $q \in Q$ in an $n$-state solvable 
	input automaton $\semiautomat{A}{Q}{\Sigma}{\delta}$
	and then test whether 
	there exists a word mapping $Q$ to $\{q\}$.
	As the latter problem is in \NP\  by Theorem~\ref{thm:solvable},
	we obtain that \SFS is in \NP. 
	\NP-hardness follows as \SFS
	is \NP-complete for binary partially ordered automata
	as proved in~\cite[Thm.~4]{Ryzhikov19a}, 
	see also Table~\ref{compl:tab:overview}.
\end{proof}

\begin{theorem} \label{thm:SYNC-solv}
\DFAST, \DFASS, \DFAS, and \SIS are \LOGSPACE-complete for solvable input automata, even for a fixed unary alphabet.
\end{theorem}
\begin{proof}
	By \cref{compl:prop:reductions} we have
	\[
	\SCP \leACZEROm \DFAST \leACZEROm  \DFASS \leACZEROm \SIS,
	\]
	and $\DFAST \leACZEROm \DFAS$,
	where the reductions of an automata-theoretical problem to another automata-theoretical problem given in the proof
	of \cref{compl:prop:reductions} preserve the alphabet.
	Since \SCP is \LOGSPACE-complete by \cref{prelim:thm:complexities},
	and the reduction to \DFAST 
	works even for unary automata, we obtain the \LOGSPACE-hardness claims even for a fixed unary alphabet.
	
	
	Since the reductions from \cref{compl:prop:reductions} transforming automata preserve
	the property of being solvable, 
	by the reductions from above, it is sufficient to prove \LOGSPACE-membership
	for \DFAS and \SIS to establish it for the other problems as well.
	
	For \DFAS, to test synchronizability, by Theorem~\ref{prelim:thm:solvrank_r} 
	we can simply test whether the word $(a_1 \cdots a_k)^{n-1}$ for $\Sigma = \{a_1, \ldots, a_k\}$
	synchronizes the $n$-state solvable input automaton, which can be done in deterministic log-space
	by Theorem~\ref{prelim:thm:finite_list_sync_word}.
	
	However, this is not sufficient since we cannot simply construct the word
	as this would take up more than logarithmic space, i.e.,
	having a deterministic log-space Turing machine 
	deciding whether a word synchronizes a given automaton,
	we cannot simply write $(a_1 \cdots a_k)^{n-1}$ and an encoding of the automaton 
	on its input tape and then simulate it, as this would
	take up more than logarithmic space.
	But we can simulate the machine by letting it read the input automaton
	from the original input tape
	and feeding
	$(a_1 \cdots a_k)^{n-1}$ one symbol at a time to it in log-space without storing the full word similarly as closure of \LOGSPACE
	is proved for log-space transducers, see~\cite[Lem.~13.2]{HopUll79}; in fact,
	a log-space transducer that additionally produces 
	$(a_1 \cdots a_k)^{n-1}$ can also be given to obtain the claim.
	% (the word \todo{satz in klammern geadded, 20.3}
	% $(a_1 \cdots a_k)^{n-1}$ cannot be stored in log-space, 
	% but it can be feed one symbol at a time to the other problem in log-space
	% similarly as closure of \LOGSPACE
	% is proved for log-space transducers, see~\cite[Lem.~13.2]{HopUll79}; in fact,
	% a log-space transducer that additionally produces 
	% $(a_1 \cdots a_k)^{n-1}$ can also be given to obtain the claim).	
		
	To solve \SIS for $\semiautomat{A}{Q}{\Sigma}{\delta}$ and $S \subseteq Q$, 
	by \crefitem{prelim:classes:solvable:thm}{prelim:prop:solvable-SCC:2}
	and \cref{prelim:classes:block:cor:Q-into-S}
we have to test whether~$S$ contains all sink components of~$\mathcal A$.
We can do this by iterating through all states $q \in Q$,
test whether it is contained in a sink component
and if this is the case, test whether $q \in S$.
By \cref{prelim:classes:block:prop:char-sink-state},
a state $q \in Q$
is in a sink component if and only if for all $x \in \Sigma$ 
we have $\delta(q, x^i) = q$ for some $i \in \{1,\ldots,|Q|\}$.
This can be checked in deterministic log-space
by successively applying~$|Q|$ times --- using a counter that counts 
up to~$|Q|$ and so takes up at most logarithmic space via a binary encoding --- the letter~$x$ to a state stored in 
a memory cell initialized with~$q$. Then it is easy to give an algorithm running
in deterministic log-space.
%and a counter that counts up to~$|Q|$ and, via a binary encoding, uses
%at most logarithmic space.	
\end{proof}

\begin{theorem} \label{prelim:compl:thm:reach-NP}
	\REACH is \NP-complete for solvable automata even for partially ordered binary automata.	
\end{theorem}
\begin{proof}
	The problem is \NP-complete for partially ordered binary
	automata~\cite[Thm.~5]{Ryzhikov19a}.
	Furthermore, as in the proof of Theorem~\ref{thm:solvable}
	by using Proposition~\ref{prop:S-T-bound}
	we obtain \NP-membership for solvable automata. 
\end{proof}

\begin{theorem} \label{prelim:compl:thm:reachf-NP}
	\REACHF is \NP-complete for solvable automata even for commutative partially ordered automata,
	and also even for partially ordered binary automata.
\end{theorem}
\begin{proof}
	\NP-hardness for partially ordered binary automata
	follows from \cref{prelim:compl:thm:reach-NP}
	and \cref{compl:prop:reductions:item:REACH-REACHF} from \cref{compl:prop:reductions}.
	Next, in~\cite[Thm.~9]{DBLP:journals/iandc/Beaudry88}
	it has been shown that deciding
	for commutative partially ordered $\mathcal A = (Q, \Sigma, \delta)$
	and $t \colon Q \to Q$
	whether there exists $u \in \Sigma^*$
	such that $\delta(q, u) = t(q)$
	for all $q \in Q$ is \NP-complete.\footnote{In~\cite{DBLP:journals/iandc/Beaudry88} this problem was formulated as a problem on transformation monoids, where the 
		commutative aperiodic transformation monoids used
		correspond precisely to the commutative partially ordered automata.}
	This problem can be reduced to \DFAINT 
	by letting $Q = \{q_1, \ldots, q_n\}$ and $\mathcal A_i = \mathcal A$
	and $F_i = \{t(q_i)\}$ for $ i \in \{1,\ldots,n\}$.
	Then when applying the reduction
	from \crefitem{compl:prop:reductions}{compl:prop:reductions:item:DFAINT_DST}
	 to obtain an instance of \ST 
	from this \DFAINT-instance
    the resulting automaton is a commutative partially ordered automaton 
	$\mathcal A' = (Q', \Sigma, \delta')$
	where $\delta'(\{q_1, \ldots, q_n\}, u) \subseteq F_1 \cup \ldots \cup F_n$
	if and only if $\delta'(\{q_1, \ldots, q_n\}, u) = F_1 \cup \ldots \cup F_n$ for $u \in \Sigma^*$, i.e., we actually have a reduction to \REACHF.
	
	As in the proof of Theorem~\ref{thm:solvable},
	by using Proposition~\ref{prop:S-T-bound},
	we obtain that \REACHF is in \NP\ 
	for solvable automata.
\end{proof}